\chapter{Evaluation}

Dieses Kapitel stellt Evaluationsdesign, Durchführung und Auswertung vor und bewertet die Anwendung anhand quantitativer Kennzahlen. Ziel ist eine reproduzierbare Gegenüberstellung der implementierten Scheduling-Verfahren unter identischen Rahmenbedingungen.

\section{Zielsetzung der Evaluation}

Die Evaluation verfolgt zwei miteinander verbundene Ziele. Erstens soll geprüft werden, ob die Implementierung fachlich plausible Unterschiede zwischen den Verfahren sichtbar macht. Zweitens soll gezeigt werden, dass die Anwendung als didaktisches Vergleichswerkzeug konsistente und interpretierbare Ergebnisse liefert.

Im Zentrum steht daher nicht die Suche nach einem universell besten Verfahren, sondern die kontextabhängige Analyse von Trade-offs zwischen Reaktionsfähigkeit, Fairness, Durchsatznähe und Warteverhalten \cite{silberschatz2018operating,tanenbaum2014modern,pinedo2016scheduling}.

\section{Evaluationsdesign}

Das Evaluationsdesign ist quantitiv und szenariobasiert. Alle Algorithmen werden auf denselben Prozessmengen ausgeführt, um Ergebnisunterschiede ausschliesslich auf das Scheduling-Verhalten zurückzuführen. Für jeden Lauf werden identische Eingangsdaten und identische Simulationsparameter verwendet.

Die Untersuchung orientiert sich an drei Szenarioklassen:

\begin{itemize}
    \item \textbf{Interaktives Lastprofil:} viele kurze Prozesse mit frühen Antworten als Ziel.
    \item \textbf{Batch-orientiertes Lastprofil:} gemischte bis lange CPU-Bursts mit Fokus auf Durchlaufzeit.
    \item \textbf{Prioritätssensitives Lastprofil:} konkurrierende Prozesse mit unterschiedlichen Prioritäten und möglichen Starvation-Effekten.
\end{itemize}

Jede Szenarioklasse wird mehrfach mit variierenden Parametern durchgeführt, um die Robustheit der Beobachtungen zu erhöhen. Die Wiederholbarkeit wird durch deterministische Laufbedingungen abgesichert, wie sie für diskrete Ereignissimulationen empfohlen werden \cite{banks2010discrete}.

\section{Durchführung der Evaluation}

Die Durchführung erfolgt in vier Schritten:

\begin{enumerate}
    \item Definition der Prozessmengen pro Szenarioklasse (Ankunftszeiten, Burst-Zeiten, Prioritäten).
    \item Ausführung aller Algorithmen pro Szenario mit festen Parametern.
    \item Erfassung der vom Simulationskern erzeugten Metriken pro Lauf.
    \item Vergleich und Interpretation der Ergebnisse innerhalb und zwischen den Szenarioklassen.
\end{enumerate}

Für Round Robin und MLFQ werden zusatzlich Parametersensitivitäten betrachtet (insbesondere Zeitquantum und Queue-Settings), da diese die Balance zwischen Reaktionszeit und Kontextwechselhäufigkeit wesentlich beeinflussen \cite{meing2025mlfq,pinedo2016scheduling}.

Die Durchführung wird protokolliert, indem pro Lauf Konfiguration und Resultate dokumentiert werden. Dadurch lassen sich die Beobachtungen auch zu einem späteren Zeitpunkt exakt nachvollziehen.

\section{Auswertung der Ergebnisse}

Die Auswertung basiert auf einem einheitlichen Metrikset:

\begin{itemize}
    \item durchschnittliche Wartezeit,
    \item durchschnittliche Durchlaufzeit,
    \item durchschnittliche Reaktionszeit,
    \item Anzahl der Präemptionen und Kontextwechsel,
    \item Fairness-bezogene Kennzahlen,
    \item Anteil inaktiver CPU-Zeit.
\end{itemize}

Die Metriken werden je Szenarioklasse tabellarisch und vergleichend dargestellt. Die Interpretation erfolgt entlang zweier Leitfragen:

\begin{itemize}
    \item Welche Verfahren sind in welcher Lastsituation erwartungsgemäss vorteilhaft?
    \item Wo entstehen Zielkonflikte zwischen kurzer Reaktionszeit, geringer Wartezeit und hoher Fairness?
\end{itemize}

Damit wird eine differenzierte Bewertung möglich, die sich an etablierten Bewertungskriterien für Scheduling orientiert \cite{silberschatz2018operating,pinedo2016scheduling}.

\section{Diskussion der Ergebnisse}

Die Diskussion ordnet die quantitativen Befunde in den didaktischen und technischen Kontext der Arbeit ein. Erwartbar ist beispielsweise, dass Verfahren mit häufigeren Präemptionen in interaktiven Szenarien bei der Reaktionszeit profitieren, gleichzeitig jedoch mehr Kontextwechsel verursachen.

Ebenso ist zu erwarten, dass prioritätsbasierte Strategien unter bestimmten Lastprofilen das Risiko unausgewogener Bedienung erhöhen, wenn keine zusätzlichen Alterungs- oder Fairnessmechanismen eingesetzt werden \cite{tanenbaum2014modern,silberschatz2018operating}.

Als Limitation ist zu berücksichtigen, dass die Evaluation auf einem abstrahierten Simulationsmodell basiert und reale Betriebssystemeffekte (z. B. Hardware-Nebenwirkungen, Interrupt-Last, komplexe I/O-Wechselwirkungen) nur eingeschränkt abbildet. Für den didaktischen Zweck der Arbeit ist diese Abstraktion jedoch gewollt, da sie den Kern der Scheduling-Entscheidungen isoliert sichtbar macht.

Insgesamt liefert die Evaluation eine nachvollziehbare Grundlage, um die Forschungsfrage evidenzbasiert zu beantworten und den Beitrag der Anwendung als interaktives Lern- und Vergleichswerkzeug zu belegen.
